\section{Sharp countermodels}

The following finite-dimensional models mark the exact limits of the preceding
theorems. Some specialize diagnostics already present in TPC-62 to the
activated carrier; they are included to close the quantifier audit, not claimed
as new counterexamples. None is an arithmetic counterexample.

\begin{example}[Ambient grouping kernel, positive activated range]
Let \(Lz=(z,z)\in\mathbb C^2\) and \(A(x_1,x_2)=x_1+x_2\). The ambient map
\(A\) has kernel \(\operatorname{span}(1,-1)\), but this kernel misses
\(Y=\operatorname{span}(1,1)\). Since
\(\|ALz\|^2/\|Lz\|^2=2\), the activated constant is \(2\). An ambient kernel
therefore does not decide the range-restricted gate.
\end{example}

\begin{example}[Every row visible, global coherent kernel]
Take \(T=I_2\) and \(F=[\,1\ -1\,]\). Each row column has squared norm \(1\),
but \(F(1,1)=0\). The normalized Gram
\[
 K=\begin{pmatrix}1&-1\\-1&1\end{pmatrix}
\]
has eigenvalues \(0,2\). Thus even perfect rowwise diagonal data do not give a
uniform represented lower bound.
\end{example}

\begin{example}[Distinguished success, uniform failure]
For \(F=[\,1\ 1\,]\), the physical-looking vector \((1,1)\) has quotient \(2\),
whereas \((1,-1)\) is killed. A favorable distinguished coefficient does not
imply a full-space generalized eigenvalue bound.
\end{example}

\begin{example}[Injectivity without a polynomial exponent]
Let \(T_X=I_2\) and \(F_X=\operatorname{diag}(1,e^{-X})\). The map is injective
for every finite \(X\), yet \(\beta_X=e^{-2X}\). Kernel exclusion alone has no
polynomial quantitative content.
\end{example}

\begin{example}[Reference route fails while a direct route succeeds]
Let terminal energy be \(T=I_2\), choose a reference projection
\(J(x_1,x_2)=x_1\), and take the final map \(A=I_2\). The comparison
\(\|Jy\|^2/\|y\|^2\) has zero infimum, while the direct \(A/T\) constant is
\(1\). Failure of a chosen intermediate reference kills that modular route,
not every direct closure.
\end{example}

\begin{example}[Stage product can be arbitrarily loose]
For \(0<\varepsilon\le1\), on \(\mathbb C^2\) take parent metric \(R=I\),
\(L=\operatorname{diag}(\sqrt\varepsilon,1)\), and
\(A=\operatorname{diag}(1,\sqrt\varepsilon)\), with no reassembly loss. The
activation and represented constants are both \(\varepsilon\), so their product
is \(\varepsilon^2\). But \(AL=\sqrt\varepsilon I\), and the exact direct
constant is \(\varepsilon\).
\end{example}

\begin{example}[Complementary singular blocks]
Let \(T=I_2\), \(F_1=[\,1\ 0\,]\), and \(F_2=[\,0\ 1\,]\). Both block Grams
are singular, but the labelled sum is \(I_2\). This saturates the kernel-fusion
criterion. If label erasure instead takes \(F_1=[\,1\ 0\,]\) and
\(F_2=[\,0\ -1\,]\) on two independent block inputs, then the coherent map is
\((x_1,x_2)\mapsto x_1-x_2\) and has the diagonal kernel. Labelled fusion and
physical fusion are distinct gates.
\end{example}

\begin{example}[Terminal-null repair can be read upstream]
Let the repaired parent be \(\mathbb C^2\), terminalization \(L(x,n)=x\), and
the original parent be the subspace \(n=0\). Every downstream map \(AL\) is
repair invariant. The pre-terminal map \(F(x,n)=x+Cn\), however, has a Gram
whose behavior changes arbitrarily with \(C\). Descent through \(L\) is a
necessary hypothesis, not notation.
\end{example}

\begin{example}[Represented success, reassembly cancellation]
Take \(A=I\), a unit vector \(v\), and \(M=-vv^*\). The represented constant is
\(1\), but \(A+M\) kills \(v\). Activation and represented transfer cannot be
used as substitutes for the reassembly short.
\end{example}

Together these examples show why the actual row Gram, its off-diagonal entries,
the label-erasure angle, and the reassembly map must be audited separately.
